\chapter{Connectivity and Strongly Connected Components}

Connectivity is the first structural question you can ask about a graph: can you get from here to there? In an undirected graph the answer carves the vertices into connected components, and a single depth-first or breadth-first sweep finds them. Directed graphs are subtler. The road network of a city is not symmetric -- a one-way street lets you go one direction but not the other -- and ``you can reach me and I can reach you'' becomes a genuine equivalence relation only when we insist on round trips. The equivalence classes of that relation are the \emph{strongly connected components} (SCCs), and they are the directed analogue of connected components.

This chapter builds the SCC machinery from the ground up. We start with the easy undirected case, define strong connectivity carefully, and then meet the two classic linear-time SCC algorithms: Kosaraju's two-pass method (transparent, easy to prove) and Tarjan's single-pass method (one DFS, the one you reach for in a contest). We prove both correct. Then we collapse each SCC to a point to obtain the \emph{condensation}, a directed acyclic graph that exposes the coarse reachability skeleton of the original digraph. Finally we cash all of this in on a real application: deciding 2-satisfiability in linear time by reading off SCCs of an implication graph.

\chapterline

\section{Undirected connectivity}
\label{ch05:undirected}

Let $\Gph = (V, E)$ be an undirected graph. Define a relation on $V$ by $u \sim v$ iff there is a path from $u$ to $v$. Reflexivity (the trivial length-0 path), symmetry (reverse the path -- legal because edges are undirected), and transitivity (concatenate paths) make $\sim$ an equivalence relation. Its classes are the \emph{connected components}.

\begin{definition}[Connected component]
\label{ch05:def-cc}
A \emph{connected component} of an undirected graph $\Gph$ is a maximal set $C \subseteq V$ such that every pair of vertices in $C$ is joined by a path. Equivalently, it is an equivalence class of the reachability relation $\sim$.
\end{definition}

\begin{intuitionbox}
Drop the graph on a table and pull on one vertex. Everything that comes up with it -- everything physically attached through a chain of edges -- is its connected component. The pieces left lying on the table are the other components. A single graph search starting at any vertex lights up exactly one component; repeat from an unvisited vertex to get the next.
\end{intuitionbox}

The algorithm is a loop over vertices that launches a DFS (or BFS) from each not-yet-seen vertex, tagging every reached vertex with a component id.

\begin{algorithm}[H]
\caption{Connected components of an undirected graph}
\label{ch05:alg-cc}
\DontPrintSemicolon
\KwIn{Undirected graph $\Gph=(V,E)$ as adjacency lists.}
\KwOut{Array $\mathrm{comp}[\cdot]$ assigning each vertex a component id.}
$c \gets 0$\;
\lForEach{$v \in V$}{$\mathrm{comp}[v] \gets \bot$}
\ForEach{$v \in V$}{
  \If{$\mathrm{comp}[v] = \bot$}{
    $c \gets c + 1$\;
    \textsc{Explore}$(v, c)$\;
  }
}
\Return $\mathrm{comp}$\;
\BlankLine
\SetKwProg{Fn}{procedure}{}{}
\Fn{\textsc{Explore}$(s, c)$}{
  push $s$ onto stack $S$; $\mathrm{comp}[s] \gets c$\;
  \While{$S$ nonempty}{
    $u \gets$ pop $S$\;
    \ForEach{$w \in \mathrm{adj}(u)$}{
      \If{$\mathrm{comp}[w] = \bot$}{
        $\mathrm{comp}[w] \gets c$; push $w$ onto $S$\;
      }
    }
  }
}
\end{algorithm}

\begin{whybox}
Every vertex is assigned a component id exactly once (the $\mathrm{comp}[w]=\bot$ guard), and the outer loop visits every vertex, so the total work is one traversal of each adjacency list: $\Oh(\abs{V}+\abs{E})$ time and $\Oh(\abs{V})$ extra space. Correctness: \textsc{Explore}$(v,c)$ assigns id $c$ to exactly the set of vertices reachable from $v$, which (by symmetry) is precisely $v$'s equivalence class, because the search follows edges in both directions and stops only when no new vertex is reachable.
\end{whybox}

\begin{connectionbox}
A union--find (disjoint set union) data structure gives the same component labels and, unlike a single static traversal, supports adding edges online: union the endpoints. That dynamic view is the entry point to the incremental connectivity and minimum-spanning-tree material later in the book. For now, a static DFS is simpler and strictly linear.
\end{connectionbox}

Undirected connectivity is the warm-up. The real content of this chapter is what happens when edges have direction.

\section{Strong connectivity: the directed story}
\label{ch05:strong}

In a directed graph $\Gph=(V,E)$, ``reach'' is one-directional. Reachability $u \rightsquigarrow v$ (a directed path from $u$ to $v$) is reflexive and transitive but \emph{not} symmetric. To recover an equivalence relation we symmetrize it deliberately.

\begin{definition}[Strong connectivity]
\label{ch05:def-scc}
Vertices $u,v$ are \emph{strongly connected}, written $u \equiv v$, if $u \rightsquigarrow v$ and $v \rightsquigarrow u$: each can reach the other. The relation $\equiv$ is an equivalence relation; its classes are the \emph{strongly connected components} (SCCs) of $\Gph$.
\end{definition}

Reflexivity and symmetry are immediate. Transitivity: if $u \rightsquigarrow v \rightsquigarrow w$ and $w \rightsquigarrow v \rightsquigarrow u$, concatenate to get $u \rightsquigarrow w$ and $w \rightsquigarrow u$. So $\equiv$ partitions $V$, and within one class every vertex can reach every other and return.

\begin{intuitionbox}
An SCC is a maximal set of vertices that all sit on a common directed cycle structure: from any one you can take a round trip through any other. If the digraph were a city of one-way streets, an SCC is a neighborhood where you can drive from any corner to any other corner and back without ever breaking the law. Between neighborhoods, travel may be one-way only.
\end{intuitionbox}

Here is a concrete digraph we will return to throughout the chapter. Read the arrows carefully.

\begin{center}
\begin{tikzpicture}[>=Stealth, node distance=16mm,
  every node/.style={circle,draw=nodeblue,fill=blue!5,minimum size=8mm,inner sep=0pt,font=\small}]
  \node (a) at (0,2) {a};
  \node (b) at (2,2) {b};
  \node (c) at (4,2) {c};
  \node (d) at (1,0) {d};
  \node (e) at (3,0) {e};
  \node (f) at (5,1) {f};
  \node (g) at (6.5,0) {g};
  \draw[->] (a) -- (b);
  \draw[->] (b) -- (c);
  \draw[->] (c) -- (a);
  \draw[->] (b) -- (d);
  \draw[->] (d) -- (e);
  \draw[->] (e) -- (b);
  \draw[->] (c) -- (f);
  \draw[->] (f) -- (g);
  \draw[->] (g) -- (f);
\end{tikzpicture}
\end{center}

Trace it by eye. The triangle $a\to b\to c\to a$ is a cycle, so $\{a,b,c\}$ are mutually reachable. But $b\to d\to e\to b$ is also a cycle, and it shares vertex $b$ with the triangle, so $a,b,c,d,e$ all reach each other: they form one SCC $\{a,b,c,d,e\}$. The pair $f\leftrightarrow g$ is a 2-cycle, hence an SCC $\{f,g\}$. Note $c\to f$ lets the big component reach $\{f,g\}$, but nothing returns from $f$ or $g$ to the rest, so the two components stay separate. Final answer: SCCs $\{a,b,c,d,e\}$ and $\{f,g\}$.

\begin{trapbox}
A common mistake is to call two vertices strongly connected just because they lie on \emph{a} directed path together. That is one-way reachability, not strong connectivity. You need a path \emph{both ways}. In the picture, $a$ reaches $g$ but $g$ cannot reach $a$, so they are in different SCCs even though a directed path $a\rightsquigarrow g$ exists.
\end{trapbox}

The naive algorithm is clear: for each pair $(u,v)$ test mutual reachability with two graph searches. That is $\Oh(\abs{V}\cdot(\abs{V}+\abs{E}))$ at best and wasteful. A slightly smarter version computes, for each vertex $v$, the set of vertices reachable from $v$ (one DFS) and the set that reach $v$ (one DFS in the reversed graph); the SCC of $v$ is the intersection. Still $\Oh(\abs{V}(\abs{V}+\abs{E}))$. We can do all of it in a single linear pass. Two algorithms achieve $\Oh(\abs{V}+\abs{E})$, and we now develop both.

\begin{historybox}
The linear-time SCC story has three landmarks. Tarjan (1972), in the same paper that systematized depth-first search, gave a single-pass algorithm using a value he called the \emph{lowlink}. Independently, the two-pass method usually attributed to Kosaraju (unpublished, 1978) and to Sharir (1981) trades one extra DFS for conceptual simplicity. Later, Cheriyan--Mehlhorn (1996) and Gabow (1999) gave a path-based single-pass algorithm with two stacks that many find the cleanest of all. We treat Kosaraju and Tarjan; Gabow appears in the exercises.
\end{historybox}

\section{Kosaraju's two-pass algorithm}
\label{ch05:kosaraju}

Kosaraju's idea is to run DFS twice: once on $\Gph$ to learn the order in which vertices \emph{finish}, and once on the reverse graph $\Gph^{\!\top}$ (every edge flipped) in decreasing finish order. The second pass, started from the right roots, peels off SCCs one at a time.

\begin{intuitionbox}
Why does reversing help? Strong connectivity is symmetric in edge direction: $u\equiv v$ in $\Gph$ exactly when $u\equiv v$ in $\Gph^{\!\top}$, because reversing every edge reverses every path but leaves ``both ways'' intact. So $\Gph$ and $\Gph^{\!\top}$ have \emph{identical} SCCs. What reversal flips is the direction of travel \emph{between} components. If in the original graph component $X$ can reach component $Y$ but not vice versa, then in the reverse graph $Y$ reaches $X$. The first DFS's finish times rank the components topologically; processing them in that ranking on the reversed graph traps each search inside a single SCC.
\end{intuitionbox}

\subsection{Finish times and the key lemma}

Recall from the DFS chapter that each vertex $v$ gets a discovery time $d[v]$ and a finish time $f[v]$, with $d[v]<f[v]$ and the parenthesis property: intervals $[d[v],f[v]]$ are properly nested or disjoint. Extend $f$ to a set of vertices by $f(C)=\max_{v\in C} f[v]$.

\begin{lemma}[Component finish order]
\label{ch05:lem-finish}
Let $C$ and $C'$ be distinct SCCs of $\Gph$. If there is an edge $(u,v)\in E$ with $u\in C$ and $v\in C'$, then $f(C) > f(C')$.
\end{lemma}

\begin{proof}
There are two cases depending on which component DFS enters first.

Case 1: $d(C) < d(C')$, where $d(C)=\min_{w\in C} d[w]$ is the first discovery time in $C$. Let $x$ be the first vertex discovered in $C$. At time $d[x]$, every vertex of $C\cup C'$ is white. Within $C$, $x$ reaches all of $C$; and through the edge $(u,v)$ it reaches $C'$ (since $u\in C$ is reachable from $x$ inside $C$, and $v\in C'$ reaches all of $C'$). By the white-path theorem, every vertex of $C\cup C'$ becomes a descendant of $x$ in the DFS forest, so they all finish before $x$ does: $f[x] = f(C) > f(C')$.

Case 2: $d(C') < d(C)$. Let $y$ be the first vertex discovered in $C'$. Could $y$ reach $C$? If it did, then via the edge $(u,v)$ (from $C$ to $C'$) we would have $C'\rightsquigarrow C\rightsquigarrow C'$, making $C$ and $C'$ mutually reachable and hence the same SCC -- contradiction. So $C'$ cannot reach $C$. Thus the entire subtree rooted at $y$ stays inside $C'$ (it cannot escape to $C$), and finishes before any vertex of $C$ is discovered. Therefore $f(C') < d(C) \le f(C)$, giving $f(C) > f(C')$ again.
\end{proof}

\begin{corollary}
\label{ch05:cor-edge}
If $(u,v)\in E$ joins two different SCCs $C\ne C'$ (with $u\in C$, $v\in C'$), then in the \emph{reverse} graph $\Gph^{\!\top}$ the edge points from $C'$ to $C$, and still $f(C)>f(C')$ where $f$ is computed on $\Gph$. So in $\Gph^{\!\top}$, edges between distinct components always run from a \emph{lower}-$f$ component to a \emph{higher}-$f$ component.
\end{corollary}

This corollary is the whole engine. Order the components by decreasing $f$. Start a DFS in $\Gph^{\!\top}$ from a vertex of the highest-$f$ component $C^\star$. Every edge leaving $C^\star$ in $\Gph^{\!\top}$ goes to a component with even higher $f$ -- but $C^\star$ already has the maximum, so there are none. The search is trapped inside $C^\star$ and recovers exactly that component. Remove it and repeat.

\subsection{The algorithm}

\begin{algorithm}[H]
\caption{Kosaraju--Sharir SCC}
\label{ch05:alg-kosaraju}
\DontPrintSemicolon
\KwIn{Digraph $\Gph=(V,E)$.}
\KwOut{Array $\mathrm{scc}[\cdot]$ labeling each vertex with its component.}
\tcp{Pass 1: DFS on $G$, push vertices on finish}
empty stack $\mathit{order}$; mark all vertices white\;
\ForEach{$v \in V$}{
  \lIf{$v$ white}{\textsc{DfsOrder}$(v)$}
}
\tcp{Pass 2: DFS on the reverse graph in decreasing finish order}
build reverse graph $\Gph^{\!\top}$; mark all vertices white; $c\gets 0$\;
\While{$\mathit{order}$ nonempty}{
  $v \gets$ pop $\mathit{order}$\;
  \If{$v$ white}{
    $c \gets c+1$\;
    \textsc{Assign}$(v, c)$ \tcp*{DFS in $G^{\!\top}$, label all reached with $c$}
  }
}
\Return $\mathrm{scc}$\;
\BlankLine
\SetKwProg{Fn}{procedure}{}{}
\Fn{\textsc{DfsOrder}$(u)$}{
  mark $u$ gray\;
  \lForEach{$w\in\mathrm{adj}_G(u)$}{\textbf{if} $w$ white \textbf{then} \textsc{DfsOrder}$(w)$}
  push $u$ onto $\mathit{order}$ \tcp*{$u$ is now finished}
}
\Fn{\textsc{Assign}$(u, c)$}{
  $\mathrm{scc}[u]\gets c$; mark $u$ gray\;
  \lForEach{$w\in\mathrm{adj}_{G^{\!\top}}(u)$}{\textbf{if} $w$ white \textbf{then} \textsc{Assign}$(w, c)$}
}
\end{algorithm}

\begin{theorem}[Correctness of Kosaraju]
\label{ch05:thm-kosaraju}
Algorithm~\ref{ch05:alg-kosaraju} labels two vertices with the same $c$ if and only if they are in the same SCC.
\end{theorem}

\begin{proof}
We show that each call \textsc{Assign}$(v,c)$ triggered from the main loop discovers exactly the SCC containing $v$. Process the components in the order the main loop launches new searches; we argue by induction that before each launch, the still-white vertices are precisely a union of whole SCCs that have not yet been output.

The main loop pops vertices in decreasing $f$ (the stack was filled in increasing finish order). Suppose we are about to launch from $v$, the highest-$f$ white vertex; let $C=\mathrm{scc}(v)$. By choice of $v$, $f(C)=f[v]$ is the maximum finish time among white vertices, so among the not-yet-output components, $C$ has the largest $f$. Now run \textsc{Assign} (a DFS in $\Gph^{\!\top}$). It certainly reaches all of $C$, since $C$ is strongly connected in $\Gph^{\!\top}$ as well. Can it escape $C$? An escaping edge in $\Gph^{\!\top}$ goes from $C$ to some other component $C'$; by Corollary~\ref{ch05:cor-edge} that requires $f(C')>f(C)$. But $C$ already has the largest $f$ among white components, and any component with larger $f$ was output earlier and is no longer white. So every escaping edge leads to a non-white, already-labeled vertex and is ignored. The search is confined to $C$, labels exactly $C$ with $c$, and the inductive hypothesis is restored. Hence the labels are exactly the SCCs.
\end{proof}

\begin{complexitybox}
Two DFS traversals plus building the reverse graph, each $\Oh(\abs{V}+\abs{E})$. Total time $\Oh(\abs{V}+\abs{E})$, extra space $\Oh(\abs{V}+\abs{E})$ for $\Gph^{\!\top}$ and the order stack. The constant factor is roughly twice a single DFS, which is why Tarjan's one-pass method is preferred when memory or constants matter.
\end{complexitybox}

\begin{labbox}
Run Kosaraju on the running digraph. Pass 1 from $a$ (adjacency in alphabetical order) finishes vertices in some order ending with $a$ last; with alphabetical scanning the finish stack (bottom first, top last) is $g,f,c,e,d,b,a$, so popping gives $a$ first. Pop $a$: \textsc{Assign} in $\Gph^{\!\top}$ from $a$ reaches $\{a,b,c,d,e\}$ (all the reverse edges of the two cycles) and labels them $c{=}1$. The original edge $c\to f$ becomes $f\to c$ in $\Gph^{\!\top}$, so from $a$ there is no way out to $f$: the search is trapped inside the first component, exactly as Corollary~\ref{ch05:cor-edge} promises. Next white top is $f$ or $g$; \textsc{Assign} from it labels $\{f,g\}$ as component 2. Two components, matching our hand analysis. Redo it choosing a different pass-1 start vertex and confirm you still get the same two SCCs.
\end{labbox}

\subsection{Tracing Kosaraju by hand}

It pays to watch the finish times and the two passes line up on the running digraph. Pass 1 is an ordinary DFS on $\Gph$ starting at $a$, scanning adjacency lists alphabetically. The discovery order is $a,b,c$ along the triangle; from $c$ the alphabetical scan reaches $f,g$ (the side branch $c\to f$) before the search returns to $b$ and descends to $d,e$. So $f$ and $g$ finish early, deep under $c$, and the two cycle vertices $d,e$ finish later. The finish times come out as below (smaller finishes first), and we push each vertex onto $\mathit{order}$ as it finishes.

\begin{center}\small
\begin{longtable}{l l l}
\toprule
Vertex & finish time $f[\cdot]$ & pushed onto $\mathit{order}$ at \\
\midrule
$g$ & $1$ & bottom \\
$f$ & $2$ & \\
$c$ & $3$ & \\
$e$ & $4$ & \\
$d$ & $5$ & \\
$b$ & $6$ & \\
$a$ & $7$ & top \\
\bottomrule
\end{longtable}
\end{center}

So $\mathit{order}$ from top to bottom is $a,b,d,e,c,f,g$ and the main loop pops $a$ first. The exact within-component finish order can vary with tie-breaking, but the \emph{component} maxima are what Lemma~\ref{ch05:lem-finish} pins down: $f(\{a,b,c,d,e\})=7$ and $f(\{f,g\})=2$, so the big component is processed first. Pass 2 runs on the reverse graph $\Gph^{\!\top}$, drawn below: every arrow of the running digraph is flipped, so in particular $c\to f$ has become $f\to c$. Starting \textsc{Assign} from $a$ can no longer escape to $\{f,g\}$.

\begin{center}
\begin{tikzpicture}[>=Stealth, node distance=16mm,
  every node/.style={circle,draw=nodeblue,fill=blue!5,minimum size=8mm,inner sep=0pt,font=\small}]
  \node (a) at (0,2) {a};
  \node (b) at (2,2) {b};
  \node (c) at (4,2) {c};
  \node (d) at (1,0) {d};
  \node (e) at (3,0) {e};
  \node (f) at (5,1) {f};
  \node (g) at (6.5,0) {g};
  \draw[->] (b) -- (a);
  \draw[->] (c) -- (b);
  \draw[->] (a) -- (c);
  \draw[->] (d) -- (b);
  \draw[->] (e) -- (d);
  \draw[->] (b) -- (e);
  \draw[->] (f) -- (c);
  \draw[->] (g) -- (f);
  \draw[->] (f) -- (g);
\end{tikzpicture}
\end{center}

The two cycles are intact (a cycle reversed is still a cycle), so $\{a,b,c,d,e\}$ and $\{f,g\}$ remain the SCCs; only the inter-component link reversed, and it now points $f\to c$ -- into the component we visit first, not out of it. Pass 2 then proceeds:

\begin{center}\small
\begin{longtable}{l l l l}
\toprule
Pop & white? & \textsc{Assign} reaches in $\Gph^{\!\top}$ & label \\
\midrule
$a$ & yes & $\{a,b,c,d,e\}$ (reverse cycle edges; $f\to c$ blocks escape) & $1$ \\
$b$ & no  & --- already labeled & --- \\
$d,e$ & no & --- & --- \\
$c$ & no  & --- & --- \\
$f$ & yes & $\{f,g\}$ (reverse 2-cycle; $f\to c$ in $\Gph$ became $c\to f$, dead in $\Gph^{\!\top}$ from here) & $2$ \\
$g$ & no  & --- & --- \\
\bottomrule
\end{longtable}
\end{center}

Two pops did real work; the rest were absorbed by an earlier \textsc{Assign}. The labels $1,2$ come out in topological order of the condensation, source first -- which is exactly why the same machinery hands you a topological sort of $\Gph^{\mathrm{SCC}}$ for nothing.

\begin{example}[A three-component chain]
\label{ch05:ex-chain}
Consider the digraph on $\{p,q,r,s\}$ with edges $p\to q$, $q\to p$, $q\to r$, $r\to s$, $s\to r$. The 2-cycles $\{p,q\}$ and $\{r,s\}$ are the only cycles, and the lone edge $q\to r$ links them one way. Kosaraju's pass 1 (start at $p$) finishes $s,r,q,p$ so $\mathit{order}$ pops $p$ first; \textsc{Assign} in $\Gph^{\!\top}$ from $p$ reaches $\{p,q\}$ (the edge $q\to r$ reversed is $r\to q$, which never fires from inside $\{p,q\}$), labels it $1$, then $r$ labels $\{r,s\}$ as $2$. The condensation is the two-vertex DAG $\{p,q\}\to\{r,s\}$. Run Tarjan on the same graph and confirm it pops $\{r,s\}$ \emph{first} (reverse topological order) and $\{p,q\}$ second -- the mirror image of Kosaraju's emission order.
\end{example}

\section{Tarjan's one-pass algorithm}
\label{ch05:tarjan}

Tarjan computes all SCCs in a single DFS by maintaining, for each vertex, a \emph{lowlink}: the smallest discovery index reachable from $v$ using tree edges followed by at most one ``back-or-cross'' edge into a vertex still on the DFS stack. When a vertex's lowlink equals its own index, it is the \emph{root} of an SCC, and that component is exactly the vertices sitting above it on an auxiliary stack.

\begin{intuitionbox}
Think of the DFS recursion as exploring a cave with a rope (the stack of vertices currently on the active path, plus everything explored under them that has not yet been packaged into a component). The \emph{index} of a vertex is the timestamp of first visit. The \emph{lowlink} answers: ``what is the oldest vertex I can loop back to, staying inside vertices that are still live on my stack?'' If a vertex can loop no further back than itself, the rope above it forms a closed pocket -- a strongly connected component -- and we cut it off.
\end{intuitionbox}

\subsection{Definitions and the algorithm}

Maintain a global counter giving each vertex a unique $\mathrm{index}[v]$ on first visit, set $\mathrm{low}[v]=\mathrm{index}[v]$ initially, and keep an explicit stack $S$ of vertices that have been discovered but not yet assigned to a finished component, with a boolean $\mathrm{onStack}[v]$. When we explore edge $(u,w)$:
\begin{itemize}
\item if $w$ is unvisited (a tree edge), recurse, then relax $\mathrm{low}[u]\gets\min(\mathrm{low}[u],\mathrm{low}[w])$;
\item if $w$ is on the stack (a back or cross edge into a live vertex), relax $\mathrm{low}[u]\gets\min(\mathrm{low}[u],\mathrm{index}[w])$;
\item if $w$ is visited but not on the stack, it already belongs to a completed SCC -- ignore it.
\end{itemize}
After exploring all of $u$'s edges, if $\mathrm{low}[u]=\mathrm{index}[u]$ then $u$ is a component root: pop $S$ down to and including $u$; those popped vertices form one SCC.

\begin{algorithm}[H]
\caption{Tarjan SCC}
\label{ch05:alg-tarjan}
\DontPrintSemicolon
\KwIn{Digraph $\Gph=(V,E)$.}
\KwOut{Array $\mathrm{scc}[\cdot]$.}
$\mathit{idx}\gets 0$; $c\gets 0$; empty stack $S$\;
\lForEach{$v\in V$}{$\mathrm{index}[v]\gets\bot$; $\mathrm{onStack}[v]\gets\textsf{false}$}
\ForEach{$v\in V$}{
  \lIf{$\mathrm{index}[v]=\bot$}{\textsc{StrongConnect}$(v)$}
}
\Return $\mathrm{scc}$\;
\BlankLine
\SetKwProg{Fn}{procedure}{}{}
\Fn{\textsc{StrongConnect}$(u)$}{
  $\mathrm{index}[u]\gets\mathit{idx}$; $\mathrm{low}[u]\gets\mathit{idx}$; $\mathit{idx}\gets\mathit{idx}+1$\;
  push $u$ onto $S$; $\mathrm{onStack}[u]\gets\textsf{true}$\;
  \ForEach{$w\in\mathrm{adj}(u)$}{
    \uIf{$\mathrm{index}[w]=\bot$}{
      \textsc{StrongConnect}$(w)$\;
      $\mathrm{low}[u]\gets\min(\mathrm{low}[u],\mathrm{low}[w])$ \tcp*{tree edge}
    }
    \uElseIf{$\mathrm{onStack}[w]$}{
      $\mathrm{low}[u]\gets\min(\mathrm{low}[u],\mathrm{index}[w])$ \tcp*{live back/cross edge}
    }
    \lElse{\textbf{ignore} \tcp*{$w$ in a finished SCC}}
  }
  \If{$\mathrm{low}[u]=\mathrm{index}[u]$}{
    $c\gets c+1$\;
    \Repeat{$w = u$}{
      $w\gets$ pop $S$; $\mathrm{onStack}[w]\gets\textsf{false}$; $\mathrm{scc}[w]\gets c$\;
    }
  }
}
\end{algorithm}

\begin{trapbox}
The single most common bug is using $\mathrm{low}[w]$ instead of $\mathrm{index}[w]$ for the back/cross-edge case (or vice versa). The tree-edge update takes $\min$ with $\mathrm{low}[w]$ because the child has finished propagating its reach; the non-tree update takes $\min$ with $\mathrm{index}[w]$, the timestamp of the target, \emph{not} its lowlink. Mixing them can merge components that should stay split, because a cross edge into a vertex $w$ does not give you $w$'s reach -- only the fact that you can touch $w$. The \texttt{onStack} test is what keeps cross edges into finished components from corrupting the lowlink: those targets are off the stack and ignored.
\end{trapbox}

\subsection{Why lowlink detects component roots}

Fix the DFS forest. For an SCC $C$, call its \emph{root} the vertex $r\in C$ with the smallest discovery index. Equivalently $r$ is the first vertex of $C$ that DFS discovers; by the white-path theorem all of $C$ lies in the subtree rooted at $r$.

\begin{lemma}[Root characterization]
\label{ch05:lem-root}
After \textsc{StrongConnect}$(u)$ finishes exploring $u$'s edges, $\mathrm{low}[u]=\mathrm{index}[u]$ if and only if $u$ is the root of its SCC.
\end{lemma}

\begin{proof}
First, $\mathrm{low}[u]$ equals the minimum index over all vertices reachable from $u$ along a path that uses tree edges down from $u$ and then at most one edge into a vertex that is still on the stack at the moment that edge is examined. A standard induction on finishing order shows $\mathrm{low}[u]=\min\{\mathrm{index}[v] : v \text{ reachable from } u \text{ by tree edges then } \le 1 \text{ live non-tree edge}\}$.

($\Leftarrow$) Suppose $u$ is the root of $C$, so $\mathrm{index}[u]$ is minimal in $C$. We claim no vertex with index $<\mathrm{index}[u]$ is reachable from $u$ by the lowlink rule, hence $\mathrm{low}[u]=\mathrm{index}[u]$. Any vertex $v$ with $\mathrm{index}[v]<\mathrm{index}[u]$ lies outside $C$ (else $u$ would not be the root). If $u$ could reach such a $v$ through the lowlink mechanism, the touched live vertex $w$ satisfies $u\rightsquigarrow w$ and, being on the stack below or around $u$, has $w\rightsquigarrow u$ (live stack vertices ancestral-or-cross-reachable establish a path back to $u$); that would put $w$ in $C$ with smaller index, contradicting minimality. So $\mathrm{low}[u]=\mathrm{index}[u]$.

($\Rightarrow$) Suppose $u$ is \emph{not} the root; let $r\ne u$ be the root of $C$, so $\mathrm{index}[r]<\mathrm{index}[u]$ and $r$ is a proper ancestor of $u$ in the DFS tree. Since $C$ is strongly connected, there is a path $u\rightsquigarrow r$. Follow it: it descends through the subtree and eventually takes a non-tree edge into an ancestor (or a cross edge into a live vertex) whose index is $<\mathrm{index}[u]$; because $r$ is still on the stack while $u$'s subtree is active, that edge is a live edge and pulls $\mathrm{low}[u]$ down to at most $\mathrm{index}[r]<\mathrm{index}[u]$. Hence $\mathrm{low}[u]<\mathrm{index}[u]$.
\end{proof}

\begin{theorem}[Correctness of Tarjan]
\label{ch05:thm-tarjan}
When \textsc{StrongConnect}$(u)$ pops the stack at a root $u$ (i.e.\ when $\mathrm{low}[u]=\mathrm{index}[u]$), the set of popped vertices is exactly the SCC of $u$.
\end{theorem}

\begin{proof}
By Lemma~\ref{ch05:lem-root} the pop happens exactly at roots. Consider the moment $u$ finishes and is found to be a root. Every vertex of $C=\mathrm{scc}(u)$ was discovered after $u$ (since $u$ is the root) and lies in $u$'s subtree, so each was pushed onto $S$ after $u$ and sits above $u$. None of them has been popped yet: a vertex is popped only when its own root finishes, and the only root in $C$ is $u$, which is finishing now. Hence all of $C$ is currently above $u$ on $S$. Conversely, any vertex $x$ strictly above $u$ on $S$ that is \emph{not} in $C$ would belong to some other SCC $C'$ whose root $r'$ lies in $u$'s subtree and has not finished -- impossible, because $r'$ would have to finish (and pop $C'$) before $u$ finishes, since $r'$ is a descendant of $u$. So the vertices above and including $u$ on $S$ are precisely $C$, and the \textbf{repeat} loop pops exactly them.
\end{proof}

\begin{complexitybox}
One DFS. Each vertex is pushed and popped once; each edge is examined once; the lowlink updates are $\Oh(1)$. Total $\Oh(\abs{V}+\abs{E})$ time, $\Oh(\abs{V})$ extra space for indices, lowlinks, the stack, and the onStack flags. No reverse graph is needed, so the constant factor is about half of Kosaraju's.
\end{complexitybox}

\subsection{Tracing Tarjan by hand}

Take a fresh small digraph designed to exercise both cycle types and a cross edge. Vertices $1,2,3,4,5$ with edges
\[
1\to 2,\quad 2\to 3,\quad 3\to 1,\quad 3\to 4,\quad 4\to 5,\quad 5\to 4,\quad 2\to 5.
\]

\begin{center}
\begin{tikzpicture}[>=Stealth, node distance=14mm,
  every node/.style={circle,draw=nodeblue,fill=blue!5,minimum size=8mm,inner sep=0pt,font=\small}]
  \node (n1) at (0,1) {1};
  \node (n2) at (2,1.6) {2};
  \node (n3) at (2,0.4) {3};
  \node (n4) at (4,1.4) {4};
  \node (n5) at (4,0.2) {5};
  \draw[->] (n1) -- (n2);
  \draw[->] (n2) -- (n3);
  \draw[->] (n3) -- (n1);
  \draw[->] (n3) -- (n4);
  \draw[->] (n4) -- (n5);
  \draw[->] (n5) -- (n4);
  \draw[->] (n2) to[bend left=10] (n5);
\end{tikzpicture}
\end{center}

We expect SCCs $\{1,2,3\}$ (the triangle) and $\{4,5\}$ (the 2-cycle); the edge $2\to 5$ is a cross edge into the $\{4,5\}$ component, and $3\to 4$ likewise reaches it one-way. Start the DFS at vertex $1$, scanning each adjacency list in increasing order. The table shows the state right after each significant event; $S$ is the Tarjan stack listed bottom-to-top.

\begin{center}\small
\begin{longtable}{l l l l l}
\toprule
Event & index & low & $S$ & note \\
\midrule
visit 1 & $\mathrm{index}[1]{=}0$ & $\mathrm{low}[1]{=}0$ & $1$ & push 1 \\
visit 2 & $\mathrm{index}[2]{=}1$ & $\mathrm{low}[2]{=}1$ & $1,2$ & tree edge $1\to2$ \\
visit 3 & $\mathrm{index}[3]{=}2$ & $\mathrm{low}[3]{=}2$ & $1,2,3$ & tree edge $2\to3$ \\
edge $3\to1$ & & $\mathrm{low}[3]{=}\min(2,0){=}0$ & $1,2,3$ & 1 on stack: use index[1]=0 \\
visit 4 & $\mathrm{index}[4]{=}3$ & $\mathrm{low}[4]{=}3$ & $1,2,3,4$ & tree edge $3\to4$ \\
visit 5 & $\mathrm{index}[5]{=}4$ & $\mathrm{low}[5]{=}4$ & $1,2,3,4,5$ & tree edge $4\to5$ \\
edge $5\to4$ & & $\mathrm{low}[5]{=}\min(4,3){=}3$ & $1,2,3,4,5$ & 4 on stack: use index[4]=3 \\
finish 5 & & $\mathrm{low}[5]{=}3\ne 4$ & $1,2,3,4,5$ & not a root \\
back at 4 & & $\mathrm{low}[4]{=}\min(3,3){=}3$ & $1,2,3,4,5$ & absorb child low \\
finish 4 & & $\mathrm{low}[4]{=}3{=}\mathrm{index}[4]$ & $1,2,3$ & \textbf{root}: pop 5,4 $\Rightarrow$ SCC$\{4,5\}$ \\
back at 3 & & $\mathrm{low}[3]{=}\min(0,3){=}0$ & $1,2,3$ & child 4 done \\
finish 3 & & $\mathrm{low}[3]{=}0\ne 2$ & $1,2,3$ & not a root \\
back at 2 & & $\mathrm{low}[2]{=}\min(1,0){=}0$ & $1,2,3$ & absorb low[3]=0 \\
edge $2\to5$ & & 5 not on stack & $1,2,3$ & \textbf{ignored} (finished SCC) \\
finish 2 & & $\mathrm{low}[2]{=}0\ne 1$ & $1,2,3$ & not a root \\
back at 1 & & $\mathrm{low}[1]{=}\min(0,0){=}0$ & $1,2,3$ & absorb low[2]=0 \\
finish 1 & & $\mathrm{low}[1]{=}0{=}\mathrm{index}[1]$ & $\varnothing$ & \textbf{root}: pop 3,2,1 $\Rightarrow$ SCC$\{1,2,3\}$ \\
\bottomrule
\end{longtable}
\end{center}

The trace highlights the two subtleties. When we examined $5\to4$, vertex $4$ was \emph{on the stack}, so we relaxed with $\mathrm{index}[4]$, dragging $\mathrm{low}[5]$ to $3$ and signaling that $5$ is not its own root. When we later examined $2\to5$, vertex $5$ had already been popped into the finished component $\{4,5\}$, so $\mathrm{onStack}[5]$ was false and the edge was correctly ignored -- exactly the guard that prevents the cross edge from merging $\{1,2,3\}$ with $\{4,5\}$.

\begin{labbox}
Re-run the same digraph but add the edge $5\to 2$. Now $5$ can reach $2$, and through the triangle $2$ reaches $5$, so all five vertices collapse into one SCC. Trace it: you will see that when finishing $4$, $\mathrm{low}[4]$ no longer equals $\mathrm{index}[4]$ (the new edge pulls a lowlink down to $1$ via $5\to 2$), so no premature root pop happens, and the single root at vertex $1$ pops all five. This is the cleanest way to feel why the lowlink test is exactly the right trigger.
\end{labbox}

\begin{trapbox}
Recursive Tarjan overflows the call stack on adversarial inputs: a path of a million vertices recurses a million deep. In production, convert \textsc{StrongConnect} to an explicit stack of (vertex, edge-iterator) frames. The logic is identical -- push a frame on a tree edge, do the lowlink relaxation when a child frame returns, run the root test when a frame's iterator is exhausted -- but you control the memory. The exercises ask you to write this version.
\end{trapbox}

\section{The condensation DAG}
\label{ch05:condensation}

Once you know the SCCs, contract each one to a single super-vertex. The result strips out all the cyclic ``noise'' and leaves the pure one-way reachability structure.

\begin{definition}[Condensation]
\label{ch05:def-condensation}
The \emph{condensation} of a digraph $\Gph$, written $\Gph^{\mathrm{SCC}}$, has one vertex per SCC of $\Gph$. There is an edge from component $X$ to component $Y$ (with $X\ne Y$) iff some edge of $\Gph$ goes from a vertex of $X$ to a vertex of $Y$. (Self-loops and parallel super-edges are dropped.)
\end{definition}

\begin{theorem}[The condensation is acyclic]
\label{ch05:thm-dag}
$\Gph^{\mathrm{SCC}}$ is a directed acyclic graph.
\end{theorem}

\begin{proof}
Suppose for contradiction $\Gph^{\mathrm{SCC}}$ contained a directed cycle $X_1\to X_2\to\cdots\to X_k\to X_1$ on distinct components. Each super-edge $X_i\to X_{i+1}$ comes from a real edge between the components, so in $\Gph$ we have $X_1\rightsquigarrow X_2\rightsquigarrow\cdots\rightsquigarrow X_k\rightsquigarrow X_1$, and within each strongly connected $X_i$ we can route freely. Then every vertex of $X_1\cup\cdots\cup X_k$ can reach every other: they all collapse into a single SCC, contradicting that the $X_i$ are distinct maximal SCCs. So no cycle exists.
\end{proof}

\begin{center}
\begin{tikzpicture}[>=Stealth, node distance=18mm]
  \begin{scope}
    \node[draw=structurepurple,fill=structurepurple!10,rounded corners,minimum width=22mm,minimum height=10mm] (C1) at (0,0) {$\{a,b,c,d,e\}$};
    \node[draw=structurepurple,fill=structurepurple!10,rounded corners,minimum width=14mm,minimum height=10mm] (C2) at (3.6,0) {$\{f,g\}$};
    \draw[->,very thick,structurepurple] (C1) -- (C2);
  \end{scope}
\end{tikzpicture}
\end{center}

For our running digraph the condensation is just two super-vertices with one edge -- the single one-way street $c\to f$ between neighborhoods. A topological order of $\Gph^{\mathrm{SCC}}$ ranks the components so that all super-edges point forward; here $\{a,b,c,d,e\}$ precedes $\{f,g\}$.

\begin{whybox}
The two SCC algorithms hand you a topological order of the condensation for free. Tarjan finishes (and pops) components in \emph{reverse} topological order of $\Gph^{\mathrm{SCC}}$: a component is popped only after all components it can reach have been popped. So numbering components $1,2,3,\dots$ in the order Tarjan emits them and then \emph{reversing} gives a valid topological sort. Kosaraju's main loop emits them in topological order directly (highest finish time first is a source of the condensation). Either way, no separate topological sort is needed.
\end{whybox}

\begin{applicationbox}
The condensation turns many ``directed'' problems into easy DAG dynamic programs. Longest path in a general digraph is NP-hard, but \emph{after} condensing, longest path in the DAG (weighting each super-vertex by, say, the number of original vertices it contains) is linear-time by DP over a topological order. Counting reachable vertices, finding a vertex that reaches everything, detecting whether the whole digraph is strongly connected (one SCC), and the 2-SAT solver below all reduce to a trivial pass over $\Gph^{\mathrm{SCC}}$.
\end{applicationbox}

\begin{connectionbox}
The condensation is the directed-graph quotient that the companion linear-algebra book would describe as collapsing along an equivalence relation. It also previews the block-triangular structure of the adjacency matrix: order vertices by SCC in reverse topological order and $\Adjm$ becomes block upper-triangular, with each diagonal block an irreducible (strongly connected) piece. That block structure is exactly what governs the convergence of Markov chains and PageRank in the spectral chapters: a chain is irreducible precisely when its transition digraph is a single SCC.
\end{connectionbox}

\section{Application: 2-SAT in linear time}
\label{ch05:2sat}

A \emph{2-SAT} instance is a Boolean formula in conjunctive normal form where every clause has at most two literals, e.g.
\[
(x_1\lor x_2)\ \land\ (\lnot x_1\lor x_3)\ \land\ (\lnot x_2\lor\lnot x_3)\ \land\ (x_3 \lor x_1).
\]
We want a truth assignment satisfying every clause, or a proof that none exists. General SAT and even 3-SAT are NP-complete, but 2-SAT is solvable in linear time, and the route is pure SCC theory.

\subsection{The implication graph}

A clause $(\alpha\lor\beta)$ is logically equivalent to the pair of implications $\lnot\alpha\Rightarrow\beta$ and $\lnot\beta\Rightarrow\alpha$ (``if $\alpha$ is false then $\beta$ must be true,'' and symmetrically). This rewriting is the key idea.

\begin{definition}[Implication graph]
\label{ch05:def-implgraph}
Build a digraph on $2n$ vertices, one for each literal $x_i$ and one for $\lnot x_i$. For every clause $(\alpha\lor\beta)$ add the two directed edges $\lnot\alpha\to\beta$ and $\lnot\beta\to\alpha$. (A unit clause $(\alpha)$ is treated as $(\alpha\lor\alpha)$, giving $\lnot\alpha\to\alpha$.)
\end{definition}

Reachability in this graph is forced truth propagation: if $u\rightsquigarrow v$ and we set $u$ true, then $v$ is forced true. The graph is \emph{skew-symmetric}: edge $\lnot\alpha\to\beta$ comes paired with $\lnot\beta\to\alpha$, so reversing every edge and swapping each literal with its negation maps the graph to itself.

\begin{center}
\begin{tikzpicture}[>=Stealth,
  every node/.style={circle,draw=nodeblue,fill=blue!5,minimum size=9mm,inner sep=1pt,font=\scriptsize}]
  \node (x1) at (0,2) {$x_1$};
  \node (nx1) at (0,0) {$\lnot x_1$};
  \node (x2) at (2.2,2) {$x_2$};
  \node (nx2) at (2.2,0) {$\lnot x_2$};
  \node (x3) at (4.4,2) {$x_3$};
  \node (nx3) at (4.4,0) {$\lnot x_3$};
  % clause (x1 v x2): ~x1->x2, ~x2->x1
  \draw[->] (nx1) to[bend left=8] (x2);
  \draw[->] (nx2) -- (x1);
  % clause (~x1 v x3): x1->x3, ~x3->~x1
  \draw[->] (x1) to[bend left=8] (x3);
  \draw[->] (nx3) to[bend right=12] (nx1);
  % clause (~x2 v ~x3): x2->~x3, x3->~x2
  \draw[->] (x2) -- (nx3);
  \draw[->] (x3) to[bend left=8] (nx2);
  % clause (x3 v x1): ~x3->x1, ~x1->x3
  \draw[->,gray] (nx3) to[bend left=22] (x1);
  \draw[->,gray] (nx1) to[bend right=22] (x3);
\end{tikzpicture}
\end{center}

\subsection{The criterion and the assignment}

\begin{theorem}[2-SAT via SCC]
\label{ch05:thm-2sat}
A 2-SAT instance is satisfiable if and only if no variable $x_i$ has $x_i$ and $\lnot x_i$ in the same SCC of the implication graph. When satisfiable, setting each literal $\ell$ to \textsf{true} iff $\mathrm{scc}(\ell)$ comes \emph{after} $\mathrm{scc}(\lnot\ell)$ in a reverse topological order of the condensation yields a valid assignment.
\end{theorem}

\begin{proof}
($\Rightarrow$, necessity) If $x_i\equiv\lnot x_i$, then $x_i\rightsquigarrow\lnot x_i$ and $\lnot x_i\rightsquigarrow x_i$. The first path forces: $x_i$ true implies $\lnot x_i$ true, a contradiction, so $x_i$ must be false. The second path forces $x_i$ true. No assignment survives both, so the instance is unsatisfiable.

($\Leftarrow$, sufficiency and construction) Assume every variable has $x_i$ and $\lnot x_i$ in different SCCs. First a symmetry fact: skew-symmetry of the implication graph means $\mathrm{scc}(\ell)$ and $\mathrm{scc}(\lnot\ell)$ are mirror images -- if $u\rightsquigarrow v$ then $\lnot v\rightsquigarrow\lnot u$, so the SCC containing $\ell$ corresponds to the SCC containing $\lnot\ell$ under literal negation, and the condensation has a matching reversal symmetry.

Compute a topological order of the condensation. Assign $\ell\gets\textsf{true}$ iff $\mathrm{scc}(\ell)$ appears later in topological order than $\mathrm{scc}(\lnot\ell)$ -- equivalently, $\mathrm{scc}(\ell)$ has the larger ``topo rank.'' Since $x_i$ and $\lnot x_i$ are in different SCCs, exactly one of them is set true, so the assignment is well-defined and consistent. We must check it satisfies every implication edge $u\to v$: we need ``$u$ true $\Rightarrow$ $v$ true,'' i.e.\ we must never have $u$ true and $v$ false.

Suppose toward contradiction $u\to v$ with $u$ true and $v$ false. ``$u$ true'' means $\mathrm{scc}(u)$ ranks after $\mathrm{scc}(\lnot u)$. ``$v$ false'' means $\mathrm{scc}(v)$ ranks before $\mathrm{scc}(\lnot v)$. The edge $u\to v$ forces $\mathrm{scc}(u)$ to rank before-or-equal $\mathrm{scc}(v)$ in topological order (edges go forward; if same component, $u\equiv v$). By skew-symmetry the mirror edge $\lnot v\to\lnot u$ forces $\mathrm{scc}(\lnot v)$ before-or-equal $\mathrm{scc}(\lnot u)$. Chain the ranks:
\[
\mathrm{rank}(\mathrm{scc}(\lnot u)) < \mathrm{rank}(\mathrm{scc}(u)) \le \mathrm{rank}(\mathrm{scc}(v)) < \mathrm{rank}(\mathrm{scc}(\lnot v)) \le \mathrm{rank}(\mathrm{scc}(\lnot u)),
\]
where the strict inequalities are ``$u$ true'' and ``$v$ false'' and the weak ones are the two implication edges. This says $\mathrm{rank}(\mathrm{scc}(\lnot u)) < \mathrm{rank}(\mathrm{scc}(\lnot u))$, a contradiction. So no edge is violated, every clause's two implications hold, and the assignment satisfies the formula.
\end{proof}

\begin{whybox}
The proof is a single clean idea: a literal is true exactly when its node ``comes later'' than its negation in the DAG of components. Skew-symmetry guarantees that ``later'' and ``earlier'' swap perfectly under negation, so the consistency you need on one edge $u\to v$ is mirrored by the consistency on $\lnot v\to\lnot u$, and the two cannot both fail.
\end{whybox}

Because Tarjan pops components in reverse topological order, the $\mathrm{scc}$ \emph{labels} it produces already satisfy ``larger reverse-topo rank = larger label.'' That makes the assignment a one-liner: set $x_i$ true iff $\mathrm{scc}(x_i) < \mathrm{scc}(\lnot x_i)$ when labels increase as Tarjan emits (reverse topological) -- or flip the comparison to match your labeling convention. Always sanity-check the direction on a tiny instance.

\begin{algorithm}[H]
\caption{2-SAT solver}
\label{ch05:alg-2sat}
\DontPrintSemicolon
\KwIn{$n$ variables, clauses each a pair of literals.}
\KwOut{A satisfying assignment, or \textsc{Unsat}.}
build implication graph $H$ on $2n$ literal vertices\;
\ForEach{clause $(\alpha\lor\beta)$}{add edges $\lnot\alpha\to\beta$ and $\lnot\beta\to\alpha$ to $H$}
run Tarjan on $H$ to get $\mathrm{scc}[\cdot]$ (labels in reverse topological order)\;
\ForEach{variable $x_i$}{
  \lIf{$\mathrm{scc}[x_i]=\mathrm{scc}[\lnot x_i]$}{\Return \textsc{Unsat}}
}
\ForEach{variable $x_i$}{
  $\mathrm{val}[x_i] \gets (\mathrm{scc}[x_i] < \mathrm{scc}[\lnot x_i])$ \tcp*{true iff its component is popped later}
}
\Return $\mathrm{val}$\;
\end{algorithm}

\begin{complexitybox}
The implication graph has $2n$ vertices and $2m$ edges for $n$ variables and $m$ clauses. Building it is $\Oh(n+m)$; Tarjan is $\Oh(n+m)$; reading off the assignment is $\Oh(n)$. Total $\Oh(n+m)$ -- linear in the size of the formula. Contrast with general SAT, where the best known algorithms are exponential in the worst case.
\end{complexitybox}

\begin{labbox}
Solve the four-clause example by hand. Build the eight directed edges, run Tarjan, and check that no $x_i$ shares a component with $\lnot x_i$. You should find a satisfying assignment such as $x_1=\textsf{true}, x_2=\textsf{false}, x_3=\textsf{true}$. Verify each clause. Then add the clause $(\lnot x_3\lor\lnot x_1)$ together with $(x_3\lor x_1)$ already present and the forcing $x_1\to x_3\to\lnot x_1$ and back; trace whether $x_1$ and $\lnot x_1$ now collapse into one SCC, which would report \textsc{Unsat}.
\end{labbox}

\begin{applicationbox}
2-SAT is the workhorse behind a surprising range of feasibility problems: placing labels on a map so none overlap (each label has two candidate positions, conflicts become clauses), scheduling with binary either/or constraints, drawing graphs with each edge routed above or below a line, and consistency checking in type systems. Whenever a problem reduces to ``each of $n$ binary choices, with pairwise constraints expressible as implications,'' the SCC method decides it in linear time.
\end{applicationbox}

\section{Connectivity refinements and edge cases}
\label{ch05:refinements}

A few structural facts round out the picture.

\textbf{A digraph is strongly connected} iff it has exactly one SCC. Testing this needs no full SCC decomposition: pick any vertex $s$; if a DFS from $s$ in $\Gph$ reaches all vertices \emph{and} a DFS from $s$ in $\Gph^{\!\top}$ reaches all vertices, then everything reaches $s$ and $s$ reaches everything, so all pairs are mutually reachable -- one SCC. Two traversals, $\Oh(\abs{V}+\abs{E})$.

\textbf{Weak connectivity}: a digraph is \emph{weakly connected} if the underlying undirected graph (forget directions) is connected. Weak connectivity is strictly easier and strictly coarser than strong connectivity -- a directed path $a\to b\to c$ is weakly connected but has three SCCs.

\textbf{Sources and sinks of the condensation}: a \emph{source} component (no incoming super-edges) and a \emph{sink} component (no outgoing super-edges) always exist in a nonempty DAG. A classic interview problem -- ``add the fewest edges to make a digraph strongly connected'' -- is answered through the condensation: if it has more than one component, the answer is $\max(\text{sources}, \text{sinks})$, where sources and sinks count source and sink super-vertices; with one component the answer is $0$.

\begin{trapbox}
Self-loops and parallel edges never change SCCs (a self-loop puts a vertex on a trivial cycle with itself, but a single vertex is already its own SCC). Isolated vertices are singleton SCCs. Make sure your implementation treats an isolated vertex as a component of size one rather than skipping it -- both Algorithm~\ref{ch05:alg-kosaraju} and Algorithm~\ref{ch05:alg-tarjan} handle this correctly because the outer loop visits every vertex, but a hand-rolled variant that only starts from in-degree-zero vertices will silently drop them.
\end{trapbox}

\section*{Summary}

Connected components fall out of one undirected traversal. Strong connectivity is the directed analogue that demands round-trip reachability; its classes, the SCCs, are found in linear time two ways. Kosaraju runs DFS on $\Gph$ to get finish times, then DFS on $\Gph^{\!\top}$ in decreasing finish order, peeling SCCs off in topological order of the condensation -- and the single lemma about component finish times proves it. Tarjan does it in one pass with the lowlink invariant, popping each SCC when its root is recognized by $\mathrm{low}[u]=\mathrm{index}[u]$, and emitting components in reverse topological order. Contracting SCCs gives the condensation, always a DAG, which exposes one-way reachability and powers DAG dynamic programs. The headline application is 2-SAT: build the implication graph, run SCC, declare unsatisfiable iff some variable shares a component with its negation, and otherwise read the assignment straight off the component order -- all in linear time.

\section{Exercises}
\label{ch05:exercises}

\begin{exercise}[Iterative Tarjan]
Rewrite \textsc{StrongConnect} (Algorithm~\ref{ch05:alg-tarjan}) without recursion, using an explicit stack of frames each holding a vertex and an index into its adjacency list. Match the recursive version exactly: perform the lowlink relaxation when a child frame returns, and the root test when a frame's adjacency list is exhausted. Argue your version is still $\Oh(\abs{V}+\abs{E})$ and uses $\Oh(\abs{V})$ explicit-stack space.
\end{exercise}

\begin{exercise}[Gabow's path-based SCC]
Gabow's algorithm replaces the lowlink number with a second stack of ``open'' vertices. On discovering $v$, push it on both the main stack $S$ and the path stack $P$. On a back/cross edge into a live vertex $w$, pop $P$ down to the vertex with $\mathrm{index}\le\mathrm{index}[w]$. When $v$ finishes and is on top of $P$, pop $S$ down to $v$ as one SCC and pop $v$ off $P$. Implement it and prove it computes the same SCCs as Tarjan. Why does $P$ encode exactly the lowlink information?
\end{exercise}

\begin{exercise}[Kosaraju needs the reverse pass]
Show by a concrete 4-vertex example that running the \emph{second} DFS on $\Gph$ (instead of $\Gph^{\!\top}$) in decreasing finish order can merge two distinct SCCs into one reported component. Pinpoint which step of the proof of Theorem~\ref{ch05:thm-kosaraju} breaks.
\end{exercise}

\begin{exercise}[Lowlink vs index]
Construct a digraph on which Tarjan, modified to relax with $\mathrm{low}[w]$ instead of $\mathrm{index}[w]$ on non-tree edges, reports the wrong SCCs. Explain in terms of cross edges why the bug appears and why it cannot appear for back edges into a direct ancestor.
\end{exercise}

\begin{exercise}[Condensation and longest path]
Given a digraph $\Gph$ where each vertex has a positive weight, give an $\Oh(\abs{V}+\abs{E})$ algorithm for the maximum total weight of vertices on any directed walk (a walk may revisit vertices). Reduce to a DAG problem on the condensation. Prove correctness and explain why allowing revisits makes the condensation the right object.
\end{exercise}

\begin{exercise}[Strong connectivity augmentation]
Prove that the minimum number of edges to add to a digraph (with more than one SCC) to make it strongly connected equals $\max(s,t)$, where $s$ and $t$ are the numbers of source and sink components in the condensation. Give a construction achieving the bound. (Hint: pair sources with sinks.)
\end{exercise}

\begin{exercise}[2-SAT assignment direction]
The text claims ``set $\ell$ true iff $\mathrm{scc}(\ell)$ comes later in topological order than $\mathrm{scc}(\lnot\ell)$.'' Tarjan emits components in reverse topological order. Determine precisely whether the correct rule with Tarjan labels is $\mathrm{scc}[\ell] < \mathrm{scc}[\lnot\ell]$ or $>$, by tracing a unit-clause instance $(x_1)$ and confirming $x_1$ comes out true.
\end{exercise}

\begin{exercise}[At-most-one constraints]
A common modeling primitive is ``at most one of $y_1,\dots,y_k$ is true.'' Naively adding a clause $(\lnot y_i\lor\lnot y_j)$ for each pair costs $\Oh(k^2)$ clauses. Design an encoding with auxiliary variables using only $\Oh(k)$ clauses that is equisatisfiable, and verify it stays within 2-SAT (every clause has $\le 2$ literals). Prove your encoding forbids exactly the assignments with two or more $y$'s true.
\end{exercise}

\begin{exercise}[Counting SCCs of a functional graph]
A \emph{functional graph} has out-degree exactly one at every vertex (each $v$ has a single successor $f(v)$). Show that every SCC of a functional graph is a single cycle, and give an $\Oh(\abs{V})$ algorithm that finds all SCCs without a general SCC routine. How does the condensation of a functional graph look?
\end{exercise}

\begin{exercise}[Reachability after condensation]
After computing the condensation $\Gph^{\mathrm{SCC}}$, you want to answer many queries ``can vertex $u$ reach vertex $v$?'' Explain how to reduce each query to reachability between two super-vertices in the DAG, and discuss the time/space tradeoff of precomputing a transitive-closure bitset over the condensation ($\Oh(\abs{V_{\mathrm{SCC}}}^2/64)$ words) versus answering each query with a fresh DAG search.
\end{exercise}

\begin{exercise}[Bridges between SCCs]
Call an edge of $\Gph$ \emph{critical} if deleting it increases the number of SCCs. Characterize critical edges in terms of the condensation and the DFS tree, and sketch a linear-time algorithm to list them. How does this relate to bridges in the undirected sense covered in a later chapter?
\end{exercise}
